\chapter{Επίλογος}

\section{Συμπεράσματα}

Η παρούσα εργασία παρουσίασε ένα σύστημα προσομοίωσης υπερτουρισμού
βασισμένο σε Μοντελοποίηση Πρακτόρων \en{(ABM)}, εφαρμοσμένο στην
πόλη της Πάτρας με χρήση πραγματικών γεωγραφικών δεδομένων 
\en{OpenStreetMap}. Το σύστημα υλοποιεί δύο τύπους πρακτόρων
(μόνιμους κατοίκους και τουρίστες) οι οποίοι αλληλεπιδρούν σε ένα
ρεαλιστικό οδικό δίκτυο, με στόχο την ποσοτικοποίηση της επίδρασης
του τουρισμού τόσο στην ποιότητα ζωής των κατοίκων όσο και στην
εμπειρία των επισκεπτών.

Τα αποτελέσματα των τριών σεναρίων προσομοίωσης αναδεικνύουν τρία βασικά
συμπεράσματα. Πρώτον, η τουριστική πίεση δεν επιδρά ομοιόμορφα στον
χώρο. Ακόμα και με δεκαπλασιασμό των τουριστών, η συντριπτική
πλειονότητα των κατοίκων διατηρεί υψηλή ευτυχία, ενώ η επίδραση
συγκεντρώνεται σε συγκεκριμένες ζώνες υψηλής τουριστικής κυκλοφορίας.
Αυτό το εύρημα υποδεικνύει ότι η χωρική διαχείριση της τουριστικής ροής
(και όχι απαραίτητα η μείωση του συνολικού αριθμού τουριστών)
μπορεί να αποτελέσει αποτελεσματικότερη στρατηγική για τη βελτίωση της
ποιότητας ζωής των κατοίκων.

Δεύτερον, η ικανοποίηση των τουριστών παρουσιάζει αξιοσημείωτη
ανθεκτικότητα. Ακόμα και στο σενάριο υψηλής πίεσης (1000 τουρίστες),
ελάχιστοι είναι οι τουρίστες που παρουσίασαν ικανοποίηση κάτω από 0.2,
ενώ τις βραδινές ώρες το 30\% είχε ανακάμψει σε τιμές 0.9–1.0.
Το εύρημα αυτό υποδηλώνει ότι η χωρική δομή της Πάτρας επιτρέπει
ικανοποιητική διασπορά της τουριστικής κίνησης για τον δοκιμασθέντα
πληθυσμό 1000 τουριστών.

Τρίτον, το σύστημα επιβεβαίωσε την αξία της ιεραρχίας \en{POI} και
του μηχανισμού \en{forward vision scan} ως εργαλείων μοντελοποίησης
ρεαλιστικής τουριστικής συμπεριφοράς. Τα αναδυόμενα μοτίβα κίνησης,
όπως συγκέντρωση γύρω από \en{Tier A POI}, παρεκτροπές σε εμπορικούς
δρόμους και βραδινή αποσυμφόρηση, αντικατοπτρίζουν παρατηρήσεις από
πραγματικά τουριστικά περιβάλλοντα και επιβεβαιώνουν την εγκυρότητα
του μοντέλου.

Συνολικά, η εργασία καταδεικνύει ότι η μοντελοποίηση με πράκτορες
αποτελεί ισχυρό εργαλείο για τη μελέτη υπερτουρισμού, επιτρέποντας
την ποσοτικοποίηση επιπτώσεων που παραδοσιακές μέθοδοι δεν μπορούν
να αποτυπώσουν. Η γεωγραφική γενικευσιμότητα του συστήματος το
καθιστά άμεσα εφαρμόσιμο και σε άλλες ελληνικές πόλεις, παρέχοντας
ένα αναπαραγώγιμο πλαίσιο ανάλυσης.

\section{Μελλοντικές Επεκτάσεις}

Το παρόν σύστημα μπορεί να επεκταθεί σε πολλαπλές κατευθύνσεις. Η
σημαντικότερη βελτίωση θα αφορούσε την ενεργή κίνηση των κατοίκων.
Aντί να είναι στατικοί, οι κάτοικοι θα μπορούσαν να κινούνται στο
οδικό δίκτυο για εργασία, αγορές και αναψυχή, αλλάζοντας δυναμικά
τη γειτνίασή τους ανάλογα με τις τουριστικές ζώνες. Αυτό θα έκανε
τη μεταβλητή ευτυχίας των κατοίκων πολύ πιο ρεαλιστική και λιγότερο
εξαρτημένη από την τυχαία αρχική τοποθέτηση τους.

Μια δεύτερη σημαντική επέκταση είναι η προσθήκη μηχανισμών πολιτικής
διαχείρισης τουρισμού, εμπνευσμένων από τη σχετική βιβλιογραφία.
Για παράδειγμα, θα μπορούσε να υλοποιηθεί σενάριο όπου οι τουρίστες
ενημερώνονται σε πραγματικό χρόνο για τον βαθμό συνωστισμού σε
\en{POI} και αποφεύγουν τα υπερφορτωμένα σημεία, ή σενάριο όπου νέα
\en{Tier A POI} προστίθενται στο μοντέλο για τη διασπορά της
τουριστικής ροής σε μεγαλύτερη έκταση.

Από τεχνική άποψη, η αντικατάσταση του \en{ContinuousSpace} του
\en{Mesa} με χωρικές δομές δεδομένων υψηλής απόδοσης
(π.χ. \en{KDtree} για ερωτήματα γειτόνων εφόσον οι πράκτορες είναι
στατικά σημεία) και η παραλληλοποίηση της εκτέλεσης θα επέτρεπαν την
κλιμάκωση σε χιλιάδες πράκτορες, επιτρέποντας τη μελέτη σεναρίων
μαζικού τουρισμού. Επιπλέον, η βαθμονόμηση του μοντέλου με εμπειρικά
δεδομένα επισκεψιμότητας, εφόσον αυτά είναι διαθέσιμα για την
εκάστοτε πόλη, θα αύξανε σημαντικά την ακρίβεια και την αξιοπιστία
των αποτελεσμάτων.

Τέλος, η ενσωμάτωση χρονικής διακύμανσης της τουριστικής ζήτησης
(εποχικότητα, εβδομαδιαία μοτίβα) και διαφορετικών προφίλ τουριστών
(\en{backpackers}, οικογένειες, ομαδικές εκδρομές) με διαφορετικές
προτιμήσεις \en{POI} και ταχύτητες κίνησης θα προσέδιδε πρόσθετο
ρεαλισμό στο μοντέλο και θα επέτρεπε τη μελέτη πιο εξειδικευμένων
σεναρίων τουριστικής διαχείρισης.

\subsubsection{Σύνδεσμος της εργασίας}
\en{https://github.com/zormpalas/ABM-overtourism}